\documentclass[11pt]{article}
\usepackage[margin=1.1in]{geometry}
\usepackage[utf8]{inputenc}
\usepackage{amsmath,amssymb}
\usepackage{microtype}
\usepackage{enumitem}
\usepackage{xcolor}
\definecolor{linkblue}{rgb}{0.1,0.25,0.55}
\usepackage[colorlinks=true,allcolors=linkblue]{hyperref}
\hypersetup{pdftitle={Mathematicians React to OpenAI's October 6, 2026 Math Release},
            pdfauthor={Lance Fortnow}}

\newcommand{\cls}[1]{\mathrm{#1}}
\setlist[itemize]{itemsep=1pt,topsep=4pt,label=$\bullet$}

\title{Mathematicians React to OpenAI's October 6, 2026 Math Release\\[4pt]
\large A review of the commentary, with links}
\author{Prepared with Claude for Lance Fortnow}
\date{October 18, 2026}

\begin{document}
\maketitle

On October 6, 2026, OpenAI posted 722 manuscripts in 372 result families to
\href{https://github.com/openai/math}{openai/math}, produced by an unreleased
internal model
(\href{https://openai.com/index/sharing-ai-progress-in-mathematics/}{OpenAI announcement}).
By OpenAI's account the model was posed roughly 4,000 problems, with an average
of about three hours of ChatGPT Pro--level compute per result. Lean
formalizations accompany some but not all results, and three manuscripts were
withdrawn within a day over a sign error
(\href{https://github.com/openai/math/blob/main/history.md}{history}). The
reactions from mathematicians fall roughly into four camps.

\section*{The awed}

Levent Alp\"oge, now at Anthropic, called it
\href{https://x.com/__alpoge__/status/2107616859595981117}{the most significant moment in mathematical history},
notable given his earlier dispute with OpenAI over its Navier--Stokes claim.
Alvaro Lozano-Robledo called October 6 possibly the most important day in the
history of mathematics, while noting that the model targeted the Riemann
Hypothesis and got only a quasi-RH. Ben Green said the release resolved roughly
three-quarters of the research goals of a European grant he had received that
June, and Roman Sauer recognized a technique he had developed over years in
solutions to two unrelated problems
(\href{https://the-decoder.com/how-much-beauty-have-we-lost-mathematicians-react-with-shock-and-disgust-as-openai-bulldozes-their-field/}{The Decoder},
summarizing the
\href{https://proofsandprompts.com/2026/10/08/100-reactions-to-100-solutions/}{Proofs and Prompts collection}
of over 100 responses).

For theoretical computer science, Scott Aaronson's
\href{https://scottaaronson.blog/?p=10169}{``The Mathocalypse''} is the
essential read. He catalogs the TCS results (the Unique Games Conjecture,
$\cls{L}=\cls{BPL}$, sub-$n\log n$ FFT and integer multiplication, the unitary
synthesis problem, parity not in $\cls{QAC}^0$, deterministic polynomial
factoring over prime fields, among others) and relays Dana Moshkovitz's
real-time reading of the UGC proof, which moved from ``unreadable'' to genuine
admiration for its strange new recursive code within a day. He also notes what
is missing: no $\cls{P}\neq\cls{NP}$, no $\cls{P}=\cls{BPP}$, no
$\cls{NEXP}\not\subseteq\cls{P/poly}$. Dakshita Khurana wrote
\href{https://dakshitakhurana.substack.com/p/classical-at-heart}{a reflective post}
on working through the AI's proof of the unitary synthesis problem she spent
years on. Constantin Kogler, who had been working on one of the problems with
OpenAI's model, described feeling like he was living in a mathematical
wonderland.

\section*{The measured}

Terence Tao's Mastodon thread introduced the ``Math 1.0 / Math 2.0'' framing
that everyone now uses
(\href{https://blog.donweb.com/terence-tao-openai-matematica-math-1-0/}{summary}):
the goal of being first to solve open problems has been optimized to the point
of unsustainability, so the field must decenter raw problem solving and value
exposition, community building and new directions. He also said the proofs
appear to contain clever new ideas that will be fruitful once digested, while
lamenting the absence of anyone who can present and teach them.

Peter Scholze urged patience---mathematics is a marathon, not a sprint---and
warned that the systems may be close to finding attacks on widely used
cryptography. Ian Agol compared the release to Perelman's 2002 proof, which
took teams years to reconstruct. Martin Bridson cautioned against letting AI
labs set the research agenda. The
\href{https://agmai.org/\#openai-release}{AGMAI statement} from the advisory
group (Gowers, Hairer, Witten and others) neither endorses nor condemns,
calling the release the beginning, not the completion, of human understanding.

On reliability, Andrew Sutherland said some proofs will likely contain
mistakes, possibly serious ones, and Alex Townsend said OpenAI should have
released the Lean-verified manuscripts first
(\href{https://rits.shanghai.nyu.edu/ai/openai-drops-722-math-manuscripts-from-an-unreleased-model}{NYU Shanghai RITS summary}).
Melissa Lee of Monash gave a balanced outsider's view in
\href{https://www.zmescience.com/science/news-science/openai-released-722-math-papers-at-once-and-mathematicians-are-freaking-out/}{The Conversation},
flagging questions about the formalizations.

\section*{The grieving}

This is the most striking part of the Proofs and Prompts collection. Hugo
Duminil-Copin said OpenAI claimed results on every major unsolved problem he
had publicly discussed in his career and that he was paralysed. Matt Zaremsky
compared four years of progress on his problem to an archaeological dig that a
trillion-dollar company blasted with TNT before handing over the skeleton. PhD
student Tristan Humbert learned that morning that the problem at the center of
his dissertation had been claimed, and spent the day rewriting his postdoc
research statement. Sam Hughes asked how much beauty had been lost; Giovanni
Mongardi described mathematicians reduced to priests of a machine god. Bryna
Kra at Northwestern worries the episode has made mathematicians more
secretive, undermining the collaboration the field depends on
(\href{https://www.tomshardware.com/tech-industry/artificial-intelligence/ai-causes-uproar-in-mathematics-as-experts-warn-against-sharing-research-with-llms-after-openai-shares-722-manuscripts-and-proofs-professor-argues-that-the-field-must-focus-on-building-community-over-raw-solutions}{Tom's Hardware}).
Danny Calegari offered a counterpoint: his curiosity-driven work was
unaffected.

\section*{The hostile}

The \href{https://www.ahmath.org/statements}{Association for Human Mathematics statement},
reposted on Tao's blog, called the release a demonstration of power rather
than scholarship and urged colleagues to stop working with OpenAI
(\href{https://mezha.ua/en/news/mathematicians-urge-end-to-openai-collaboration-315962/}{coverage}).
Its line that mathematicians did not ask for this work to be done drew
mockery, including from Aaronson, who compared it to Hardy asking Ramanujan
who requested his identities. Henry Wilton said this would be a great way to
destroy the mathematical community and criticized the absence of named authors
and failure rates. Enrico Fatighenti objected that AI papers get a pass on
writing quality that would sink a human submission, while the community bears
the cost of deciphering them. Michael Harris at Columbia said he will not read
the papers, and David Madore offered the
\href{https://x.com/gro_tsen/status/2108217461577294035}{crudest summary}.

Straight news coverage:
\href{https://www.nature.com/articles/d41586-026-03196-8}{Nature},
\href{https://www.scientificamerican.com/article/openai-unleashes-hundreds-more-math-results-upon-a-field-already-in-shock/}{Scientific American},
\href{https://futurism.com/artificial-intelligence/mathematicians-openai-furious}{Futurism}.

\section*{Assessment}

The durable criticism is about process, not entitlement. Withholding the model
name, the prompts and the failure rates, and dumping undigested proofs,
deserve the pushback they are getting. The AHM's ``nobody asked'' position, by
contrast, will not survive the model's public release. Aaronson's contrast
between the ``OpenAI model'' (dump raw proofs and let humans race to digest
them) and the ``Anthropic model'' (give Vassilevska Williams and Alman time to
write up the 3SUM/APSP result) may prove more consequential than any single
theorem, since both leave a company deciding who gets credit, just in
different ways. Neither addresses Tao's real worry: what happens to the
fertility of a field once its open problems have been strip-mined.

\section*{Sources}

\begin{itemize}
\item OpenAI,
  \href{https://openai.com/index/sharing-ai-progress-in-mathematics/}{Sharing AI progress in mathematics}
  and \href{https://github.com/openai/math}{openai/math}
\item The Decoder,
  \href{https://the-decoder.com/how-much-beauty-have-we-lost-mathematicians-react-with-shock-and-disgust-as-openai-bulldozes-their-field/}{``How much beauty have we lost?''}
\item Proofs and Prompts,
  \href{https://proofsandprompts.com/2026/10/08/100-reactions-to-100-solutions/}{100 reactions to 100 solutions}
\item Scott Aaronson, \href{https://scottaaronson.blog/?p=10169}{The Mathocalypse}
\item Dakshita Khurana,
  \href{https://dakshitakhurana.substack.com/p/classical-at-heart}{Classical at heart}
\item AGMAI, \href{https://agmai.org/\#openai-release}{Statement on the OpenAI release}
\item Association for Human Mathematics,
  \href{https://www.ahmath.org/statements}{Statements}
\item Melissa Lee,
  \href{https://www.zmescience.com/science/news-science/openai-released-722-math-papers-at-once-and-mathematicians-are-freaking-out/}{The Conversation, via ZME Science}
\item \href{https://www.nature.com/articles/d41586-026-03196-8}{Nature},
  \href{https://www.scientificamerican.com/article/openai-unleashes-hundreds-more-math-results-upon-a-field-already-in-shock/}{Scientific American},
  \href{https://www.tomshardware.com/tech-industry/artificial-intelligence/ai-causes-uproar-in-mathematics-as-experts-warn-against-sharing-research-with-llms-after-openai-shares-722-manuscripts-and-proofs-professor-argues-that-the-field-must-focus-on-building-community-over-raw-solutions}{Tom's Hardware},
  \href{https://futurism.com/artificial-intelligence/mathematicians-openai-furious}{Futurism}
\item \href{https://rits.shanghai.nyu.edu/ai/openai-drops-722-math-manuscripts-from-an-unreleased-model}{NYU Shanghai RITS summary}
\item \href{https://blog.donweb.com/terence-tao-openai-matematica-math-1-0/}{Summary of Tao's Mastodon thread}
\end{itemize}

\end{document}
